\documentclass[11pt]{article}
\usepackage[margin=1in]{geometry}
\usepackage{times}
\usepackage{hyperref}
\hypersetup{colorlinks=true, linkcolor=black, urlcolor=blue, citecolor=black}
\title{The Patent Office Was Right the First Time: Joe Newman's Energy Machine}
\author{Flyxion \\ \small Independent Researcher}
\date{}
\begin{document}
\maketitle

\begin{abstract}
Joe Newman's Energy Machine, a large rotating assembly of permanent magnets and coils patented and promoted from the late 1970s onward, was claimed to produce more usable energy output than electrical input, a claim tested directly by the National Bureau of Standards in 1986 at a federal court's request and found, under conditions Newman himself helped design, to produce no more energy than it consumed once all output was properly measured. This paper treats the case as unusual among the claims surveyed in this series precisely because it received a rare, formal, court-ordered independent test with a clear negative result, and examines why the claim nonetheless continued to circulate for decades afterward among a persistent community of supporters.
\end{abstract}

\section{Introduction}

Joe Newman claimed from the late 1970s that his Energy Machine, a large assembly of permanent magnets and wound coils driven in rotation, produced substantially more usable energy than the electrical energy supplied to it, in violation of standard electromagnetic theory. Newman sued the United States Patent and Trademark Office after his patent application was rejected, and the resulting federal court case led to a rare event in the history of these claims: a formal, independently conducted efficiency test performed by the National Bureau of Standards, now the National Institute of Standards and Technology, in 1986, with Newman's own participation in the test's design and setup.

\section{What the Test Actually Found}

The National Bureau of Standards test measured both electrical input and mechanical and electrical output from Newman's device under conditions negotiated with Newman's own involvement, specifically to address his prior objections that previous informal tests had measured output incorrectly. The measured result showed the device's output did not exceed its input once all forms of output were properly accounted for, consistent with ordinary electromagnetic theory and inconsistent with Newman's over-unity claim. This is among the more decisive negative results in the history of these claims specifically because it was conducted by a neutral government metrology body under court order, with the inventor's own input into the test protocol, removing much of the "the test wasn't fair" objection available to less formally tested devices elsewhere in this genre.

\section{Why the Claim Continued to Circulate Anyway}

Despite this result, Newman and a body of supporters continued promoting the device for years afterward, generally by disputing specific measurement methodology details of the 1986 test, arguing that the device's claimed benefit lay in a form of energy or a measurement quantity, sometimes described in terms of pulsed rather than continuous power delivery, that the Bureau's instrumentation was not configured to properly capture. This response pattern, disputing the specific instrumentation of a negative test after having helped design that instrumentation, rather than arranging or requesting a follow-up test addressing the specific objection, recurs across the over-unity genre and is functionally difficult to distinguish from simply declining to accept a result inconvenient to the underlying claim, regardless of the objection's technical merit in isolation.

\section{Why This Case Is Instructive for the Rest of the Series}

Most devices surveyed elsewhere in this series have never received a test as formal, as independently conducted, or as directly responsive to the inventor's own stated concerns as the Newman device did in 1986, which makes the continued circulation of the claim afterward a useful data point: even a rare, well-designed, court-ordered, negative independent test does not, on its own, end a claim's circulation once a community of supporters has formed around it, suggesting that the persistent absence of independent testing common to most other devices in this series is not the sole obstacle to public skepticism resolving these claims, and that the underlying appeal of an inventor persecuted by an indifferent scientific establishment operates somewhat independently of the specific evidentiary record available.

\section{The Entropic Gravity Framing}

Newman's claim, like the Yildiz Motor and other over-unity devices addressed elsewhere in this series, concerns energy conservation rather than gravitational coupling directly, and no gravitational claim attaches to it, making it relevant to this series primarily as a comparative case study in how independent testing, even when unusually rigorous and even when directly responsive to an inventor's own objections, interacts with the broader over-unity and fringe-energy claim ecosystem this series otherwise treats.

\section{Objections}

Supporters have argued in the years since that Newman's device anticipated later, more sophisticated pulsed-motor and coil-switching designs that achieved genuine efficiency improvements over simple DC motors, and that the 1986 test's negative finding does not preclude the existence of real, if more modest, efficiency gains from asymmetric pulsed coil driving unrelated to the specific over-unity claim originally made. Efficiency improvements in motor design through better switching and pulse timing are a genuine area of conventional electrical engineering, but a genuine efficiency improvement within the bounds of energy conservation is a categorically different claim from the over-unity claim the 1986 test specifically addressed and found unsupported, and conflating the two does not rehabilitate the original claim.

\section{Conclusion}

The Newman Energy Machine received one of the more formally conducted and directly responsive independent efficiency tests in the history of over-unity claims, found no anomalous energy output, and continued to circulate among supporters afterward largely unaffected by that result, a pattern that suggests the persistence of these claims owes as much to their appeal as narratives of persecuted inventors as to any genuine ambiguity remaining in their evidentiary record.

\end{document}
